\documentclass[10pt,a4paper]{amsart}
\usepackage[margin=19mm]{geometry}
\usepackage{amsmath,amssymb,mathtools}
\usepackage[scr=rsfs,cal=euler]{mathalfa}
\usepackage{xfrac}
\usepackage[unicode,hidelinks,bookmarks=false]{hyperref}
\newcommand{\ie}{i.e.\ }
\newcommand{\cf}{cf.\ }
\newcommand{\resp}{respectively\ }
\newcommand{\Subsection}{\subsection}
\providecommand{\nobreakdash}{\nobreak\hbox{-}}
\setlength{\emergencystretch}{2em}
\multlinegap0pt
\usepackage{amssymb}
\newtheorem{lemma}{Lemma}
\newcommand{\D}{\mathcal{D}}
\newcommand{\NN}{\mathbb{N}}
\newcommand{\RR}{\mathbb{R}}
\newcommand{\Ss}{\mathcal{S}}
\newcommand{\mbf}{\mathbf} 
\title{Singularity of Data Analytic Operations}
\author{Steven P. Ellis}
\date{Source: arXiv:1307.7624v8 (CC BY 4.0)}
\begin{document}
\maketitle

\noindent\emph{Source and licence.} Singularity of Data Analytic Operations. Version arXiv:1307.7624v8 (CC BY 4.0), distributed by arXiv under the Creative Commons Attribution 4.0 International licence, http://creativecommons.org/licenses/by/4.0/, as recorded in the arXiv licence metadata for this exact version and shown on the abstract page https://arxiv.org/abs/1307.7624v8. Full licence notice: \texttt{licenses/arxiv-1307.7624-CC-BY-4.0.txt}.\par\smallskip

\noindent\textit{Editorial scope.} Counted unit: the article's lemma L:D'.is.dense.for.aug.median (source0.tex lines 12617--12619). The article's complete original proof of it is delivered with it (source0.tex lines 19854--20093, one \texttt{proof} environment of 240 lines, which the article places away from the statement). No mathematics is written, repaired or re-derived: the statement and the proof are the article's own lines, in the article's own notation, and no retained line is substituted or re-flowed.  Retained uncounted as the source-local context the proof needs: lines 697--700; lines 12415--12418; lines 12484--12488; lines 12493--12498; lines 12508--12512; lines 12523--12525; lines 12603--12606.  Nothing was omitted from any retained span, so the package carries no omission record.  No retained line contains a citation, so the wrapper carries no bibliography and no bibliography entry.  No retained line contains a figure environment, an image-inclusion command or drawing code, so no image is delivered and the package declares no asset dependency.  Editorial formatting only, in addition: an AMS \texttt{amsart} preamble that restates the article's own declarations and macros used by the retained lines, namely the environments \texttt{lemma} on separate counters, together with the standard packages those lines need.  The retained environments print in this document as Lemma 1 in order of appearance, whereas the article prints the same items with the numbering of its own full document.  The complete native source of the article as distributed in the arXiv e-print for this exact version is bundled unchanged as \texttt{sources/100087/source0.tex}.
   \begin{equation}  \label{E:D'.=.D.less.S}
      \D' = \D \setminus \Ss, \; \Phi \text{ is continuous on } \D', 
        \text{ and $\D'$ is dense in $\D$}.
   \end{equation} 

	\begin{equation}  \label{E:Ga.defn}
		G_{a}(v; x) := G_{a}(v) := a \angle(v, y_{0}) +  \sum_{i=1}^{n} \angle (y_{i}, v),  
		            \quad v \in S^{1},
	\end{equation}

	\begin{multline}  \label{E:angle.deriv.sign}
		\text{If } y \in S^{1} \setminus \{ v, -v \} \text{ then } 
		  \frac{d}{dt} \angle \bigl( y,\phi(t) \bigr) \restriction_{t=0} \text{ exists and } \\
		    \frac{d}{dt} \angle \bigl( y,\phi(t) \bigr) \restriction_{t=0} = -sign(y \cdot w).
	\end{multline}  

    \begin{align} \label{E:one-sided.angle.derivs}
      \frac{d_{+}}{dt} \angle \bigl( v,\phi(t) \bigr) \restriction_{t=0} 
        &= \frac{d_{-}}{dt} \angle \bigl( -v,\phi(t) \bigr) \restriction_{t=0} = 1 \text{ and } \\
      \frac{d_{-}}{dt} \angle \bigl( v,\phi(t) \bigr) \restriction_{t=0} 
        &= \frac{d_{+}}{dt} \angle \bigl( -v,\phi(t) \bigr) \restriction_{t=0} = -1 . \notag
    \end{align}

	\begin{multline} \label{E:Y(x).defn}
	     (y_{1}, \ldots, y _{n}) := x, \;
		  Y := Y(x)  := \{ y_{0}, y_{1}, \ldots, y_{n} \} \subset S^{1} \\
		           \text{ and } - Y := -Y(x) := \{ -y_{0}, -y_{1}, \ldots, -y_{n} \} \subset S^{1}.
	\end{multline}

	\begin{equation} \label{E:Ma(x).in.Y}
		M_{a}(x) \subset Y(x).
	\end{equation}

  \begin{lemma} \label{L:not.both.antipodals.in.Ma}
If $z_{\beta} = -z_{\alpha}$ ($\alpha, \beta = 0, \ldots, t$) then at most 
one of $z_{\alpha} $ and $z_{\beta}$ is in $M_{a}(x)$.
  \end{lemma}

  \begin{lemma} \label{L:D'.is.dense.for.aug.median}
$\D'$ is dense in $\D$ and every point in $(\D')^{c} := \D \setminus \D'$ is a singularity of $m_{a}$. Thus, $(\Phi, \D')$ with $\Phi = m_{a}$ satisfies \eqref{E:D'.=.D.less.S}.
  \end{lemma}

  \begin{proof}[Proof of lemma \ref{L:D'.is.dense.for.aug.median}]
Let $z_{1}, \ldots, z_{t} \in S^{1}$ be the distinct locations of the points 
$y_{1}, \ldots, y_{n}$. We only count observations, i.e.\ $y_{1}, \ldots, y_{n}$, not the augmentation point $y_{0}$, but it is possible that 
$y_{0} \in \{ z_{1}, \ldots, z_{t} \}$. 
Let $\ell_{\alpha} \in [1,n]$ be the multiplicity 
of $z_{\alpha}$ ($\alpha=1, \ldots, t$) 
so $\ell_{1} + \cdots + \ell_{t} = n$. Call $\{ z_{1}, \ldots, z_{t} \}$ the ``support'' of $x$. 
If for every $j \in \NN_{n}$ we have $y_{j} \neq y_{0}$, define $z_{0} = y_{0}$. 

Suppose $x = ( y_{1}, \ldots, y_{n} ) \in \D \setminus \D'$, 
then there exist $z_{i}, z_{j} \in M_{a}(x) \subset Y(x)  := \{ y_{0}, y_{1}, \ldots, y_{n} \}$,
but $z_{i} \neq z_{j}$. Let 
    \begin{equation*}
      g := g(x) := \min_{v \in S^{1}} G_{a}(v;x) . 
    \end{equation*}
Choose $z_{i}, z_{j} \in M_{a}(x)$ to have minimal separation in terms of arc length of any pair of points in $M_{a}(x)$. Thus, $t > 1$ and $G_{a}(z_{j};x) = G_{a}(z_{i};x) = g$. Either $i$ or $j$ may be 0. WLOG $j = 1$. $z_{i} \neq z_{1}$ and, by lemma \ref{L:not.both.antipodals.in.Ma}, $z_{i} \neq -z_{1}$. We may always take 
$z_{1} \neq y_{0}$ because if $z_{1} = y_{0}$ we may simply swap the indices $i$ and 1.

Let $u, v \in S^{1}$ with $v \notin \{u,-u\}$. The two arcs in $S^{1}$ joining $u$ and $v$ have positive length and, since $v \neq u$ and $v \neq -u$, one arc has length $< \pi$, the other $> \pi$. Say that a point of $S^{1}$ lies ``between $u$ and $v$'' if it lies in the shorter arc. Thus, $u$ and $v$ lie between $u$ and $v$. A point lying in the interior of the shorter arc lies ``strictly between'' $u$ and $v$ so neither $u$ nor $v$ lie strictly between $u$ and $v$

Let 
$B := \bigl\{ z \in Y(x) \cup (-Y(x)) : z \text{ lies strictly between } z_{i} \text{ and } z_{1} \bigr\}$. (See \eqref{E:Y(x).defn}.) \emph{Claim:} $B \neq \varnothing$. 
For suppose $B = \varnothing$. Then obviously no point of $Y(x)$ or $-Y(x)$ lies strictly between $z_{i}$ and $z_{1}$. If $z \in Y(x) \cup (-Y(x))$ lay strictly between $-z_{i}$ and $-z_{1}$ then $-z \in Y(x) \cup (-Y(x))$ and lies between $z_{i}$ and $z_{1}$ so that is also ruled out. We conclude that 
$B = \varnothing$ implies any $z \in Y(x) \cup (-Y(x))$ lies between $z_{1}$ and $-z_{i}$ or between $z_{i}$ and $-z_{1}$. (Draw a picture.)

Continue to suppose $B = \varnothing$. WLOG $z_{1} = (1,0)$. Measure angles counterclockwise from $z_{1}$. Then, since $z_{i} \neq - z_{1}$, we may assume 
$\theta := \angle(z_{i}, z_{1}) \in (0, \pi)$. Let $v = (\cos \phi, \sin \phi) \in S^{1}$ lie strictly between $z_{i}$ and $z_{1}$. Thus, we may take $\phi \in (0, \theta)$. 
Take $w := w(v) := \bigl( \cos (\phi+\pi/2), \sin (\phi+\pi/2) \bigr) \in S^{1}$, 
so $w \perp v$. 
If $z = (\cos \zeta, \sin \zeta)  \in S^{1}$, then 
$z \cdot w = \cos (\zeta - \phi - \pi/2) = -\sin(\zeta - \phi)$. Suppose $z$ lies between $z_{1}$ and $-z_{i}$. We may take $\zeta \in [-\theta, 0]$. I.e., both $\zeta$ and $\phi$ lie in the interval $[-\theta, \theta]$. Hence, $\zeta - \phi \in [-\pi, 0]$ and 
$z \cdot w = -\sin(\zeta - \phi) \in [0,1]$. I.e., for $z$ between $z_{1}$ and $-z_{i}$, the sign of $z \cdot w(v)$ is constant in $v$ strictly between $z_{i}$ and $z_{1}$. 
Similarly, for $z$ between $-z_{1}$ and $z_{i}$ the sign of $z \cdot w(v)$ is constant in  in $v$ strictly between $z_{i}$ and $z_{1}$. To sum up, if $B = \varnothing$ then for every $z \in Y(x) \cup (-Y(x))$, the sign of $z \cdot w(v)$ is constant in  in $v$ strictly between $z_{i}$ and $z_{1}$. 

It follows from \eqref{E:angle.deriv.sign} that if $B = \varnothing$ then the derivative 
of $G_{a}(v; x)$ is constant in $v$ between $z_{i}$ and $z_{1}$. But $v=z_{i}$ and $v=z_{1}$ are both minimizers of $G_{a}(v; x)$. Therefore, the derivative of $G_{a}(v; x)$ between $z_{i}$ and $z_{1}$ is 0, i.e., every point between between $z_{i}$ and $z_{1}$ minimizes $G_{a}(\cdot; x)$. This contradicts \eqref{E:Ma(x).in.Y} and the claim that $B \neq \varnothing$ is proved. 
 
Let $z \in B$. So $z \notin \{ z_{1}, z_{i} \}$. Since $z_{i}, z_{j}$ have minimal separation in terms of arc length of any pair of points in $M_{a}(x)$, we have 
$z \notin M_{a}(x)$. There exists $\alpha = 1, \ldots, t$ s.t.\ $z = \pm z_{\alpha}$.  
Let $\phi : (-\pi, \pi] \to S^{1}$ parametrize $S^{1}$ by arc length 
with $\phi(0) = z_{\alpha}$ and $\phi^{-1}(z_{1}) < 0$. So $\phi^{-1}(z_{i}) > 0$. 
Let $\beta = 0, 1, \ldots, t$. If $z_{\beta} = \phi(s)$ with $s \in (-\pi, 0)$ say that $\beta$ and $z_{\beta}$ are ``early''. For example, $\beta = 1$ is early. If $z_{\beta} = \phi(s)$ with $s \in (0,\pi)$ say that $\beta$ is ``late''. For example, $\beta = i$ is late. Let $E(\beta) := 1$ or 0 according as $\beta$ is early or late. 
Thus, if $\beta = 0, \ldots, t$ is neither early nor late, then we must have $z_{\beta} = \pm z_{\alpha}$. If $-z_{\alpha} \in \{ z_{1}, \ldots, z_{t} \}$ 
let $\ell_{-\alpha}$ be the corresponding weight. 
If $-z_{\alpha} \notin \{ z_{1}, \ldots, z_{t} \}$, let $\ell_{-\alpha} := 0$. 
  
For $s \in \RR$ close to 0, define $x(s) \subset S^{1}$ to be the data set with support 
$\{ z_{1}(s), \ldots, z_{t}(s) \}$, where $z_{\beta}(s) = z_{\beta}$
if $\beta \neq \alpha$, and $z_{\alpha}(s) = \phi(s)$. We assume $|s|$ is small enough that 
$\{ z_{1}(s), \ldots, z_{t}(s) \}$ are distinct. Assign to $z_{\beta}(s)$ the same multiplicity 
$\ell_{\beta}$ as before. 

We analyze $M_{a} \bigl( x(s) \bigr)$ for small $|s|$. 
By \eqref{E:Ga.defn}, $G_{a}(v; x)$ is continuous in $x$. Therefore, 
if $z_{\beta} \notin M_{a}(x)$ then for $|s|$ sufficiently small 
$z_{\beta}(s) \notin M_{a} \bigl( x(s) \bigr)$. 

Let $z_{\beta} \in M_{a}(x) \setminus \{ \pm z_{\alpha} \}$. ($z \notin M_{a}(x)$ but $\pm z_{\alpha}$ might equal $-z$ and it is possible that $-z \in M_{a}(x)$.)
Then $z_{\beta}(s) = z_{\beta}$. Notice that 
$\angle \bigl(  z_{\beta} , z_{\alpha}(s) \bigr) = \angle \bigl(  z_{\beta} , z_{\alpha} \bigr) + \bigl( 2 E(\beta) - 1 \bigr) s$. Then
    \begin{align}   \label{E:Ga(z.beta(s),x(s)}
      G_{a} \bigl( z_{\beta}(s); x(s) \bigr) 
        &= G_{a} \bigl( z_{\beta}; x(s) \bigr) \notag \\
        &= a \angle(z_{\beta}, y_{0}) 
          + \ell_{\alpha} \angle( z_{\beta}, z_{\alpha}(s) )  \notag \\
        & \qquad + \ell_{-\alpha} \angle(z_{\beta}, -z_{\alpha}) 
          + \sum_{\gamma \text{ early}} 
            \ell_{\gamma} \angle( z_{\gamma}, z_{\beta} ) 
                + \sum_{\gamma \text{ late}} \ell_{\gamma} 
                  \angle( z_{\gamma}, z_{\beta} ) \notag \\
        &= a \angle(z_{\beta}, y_{0}) + \ell_{\alpha} \angle( z_{\beta}, z_{\alpha} )
           +  \ell_{\alpha} \bigl( 2 E(\beta) - 1 \bigr) s  \\
                   & \qquad + \ell_{-\alpha} \angle(z_{\beta}, -z_{\alpha}) 
          + \sum_{\gamma \text{ early}} 
            \ell_{\gamma} \angle( z_{\gamma}, z_{\beta} ) 
                + \sum_{\gamma \text{ late}} \ell_{\gamma} 
                  \angle( z_{\gamma}, z_{\beta} ) \notag \\
        &= G_{a}(z_{\beta}; x) + \ell_{\alpha} \bigl( 2 E(\beta) - 1 \bigr) s \notag \\
        &= g + \ell_{\alpha} \bigl( 2 E(\beta) - 1 \bigr) s , \notag
    \end{align}
since $z_{\beta} \in M_{a}(x)$. This implies the following. Suppose \emph{$s$ is negative} but close to 0. Then, first, if $\beta$ is early then 
$z_{\beta} \in M_{a} \bigl( x(s) \bigr)$ and 
$G_{a} \bigl( z_{\beta}(s); x(s) \bigr) < g$. And second, if $\beta$ is late then 
$z_{\beta} \notin M_{a} \bigl( x(s) \bigr)$ and 
$G_{a} \bigl( z_{\beta}(s); x(s) \bigr) > g$. 

If \emph{$s$ is positive} but close to 0 we have the opposite: If $\beta$ is late then $z_{\beta} \in M_{a} \bigl( x(s) \bigr)$ and 
$G_{a} \bigl( z_{\beta}(s); x(s) \bigr) < g$. If $\beta$ is early then 
$z_{\beta} \notin M_{a} \bigl( x(s) \bigr)$ and 
$G_{a} \bigl( z_{\beta}(s); x(s) \bigr) > g$. 

Recall that $z = \pm z_{\alpha}$ is a point in $B$. First assume 
    \begin{equation*}
      \text{Neither $z$ nor $-z$ is in } M_{a}(x). 
    \end{equation*}
Then every point of $M_{a}(x)$ is either early or late and we have
    \begin{align} \label{E:precisely.in.Ma(x(s))}
      &\text{If $s$ is negative but close to 0 then the points of  } \notag
        M_{a} \bigl( x(s) \bigr) \text{ are precisely } \\ 
      & \qquad \text{ the early points of $M_{a}(x)$ and } \\
      &\text{If $s$ is positive but close to 0 then the points of  } \notag
        M_{a} \bigl( x(s) \bigr) \text{ are precisely } \\ \notag
      & \qquad \text{ the late points of $M_{a}(x)$. } \notag
    \end{align}

Now assume 
    \begin{equation*}
      z_{\alpha} \in M_{a}(x).
    \end{equation*}
Since $z \notin M_{a}(x)$ we must have $z_{\alpha} = -z$. We show that 
$z_{\alpha}(s) \notin M_{a} \bigl( x(s) \bigr)$ so 
\eqref{E:precisely.in.Ma(x(s))} \emph{continues to hold}.

Consider $G_{a} \bigl( z_{\alpha}(s); x(s) \bigr)$. 
Since $z \notin M_{a}(x)$, we must have $z_{\alpha} = -z \notin B$, 
so $z \neq z_{\alpha}$ and $z_{\alpha} = \phi(\pi)$. 
To compute $G_{a} \bigl( z_{\alpha}(s); x(s) \bigr)$ we need to compute 
$\angle \bigl( y_{0}, z_{\alpha}(s) \bigr)$ for $|s|$ small:
    \begin{equation}  \label{E:angle(y0,z.alpha(s))} 
      \angle \bigl( y_{0}, z_{\alpha}(s) \bigr) =
        \begin{cases}
             \angle ( y_{0}, z_{\alpha} ) + \bigl( 2 E(0) - 1 \bigr) s, 
               &\text{ if } z_{\alpha} \notin \{ \pm y_{0} \}, \\
             |s| = \angle ( y_{0}, z_{\alpha} ) + |s|, &\text{ if } y_{0} = z_{\alpha}, \\
             \pi - |s| = \angle ( y_{0}, z_{\alpha} ) - |s|, &\text{ if } y_{0} = -z_{\alpha} .
        \end{cases}
    \end{equation}
We summarize this by writing 
$\angle \bigl( y_{0}, z_{\alpha}(s) \bigr) = \angle ( y_{0}, z_{\alpha} ) + \Delta_{0} s$. 
So $\Delta_{0} = \Delta_{0}(s) = \Delta_{0}(s;y_{0}) = \pm 1$. Note also that 
$\angle \bigl( z_{\alpha}(s), -z_{\alpha} \bigr) = \angle(z_{\alpha}, -z_{\alpha}) - |s|$. Similarly, 
$\angle \bigl( z_{\alpha}(s), z_{\alpha}(s) ) =  \angle(z_{\alpha}, z_{\alpha}) = 0$. We have then,
    \begin{align}   \label{E:Ga(z.alpha(s),x(s)}
      G_{a} \bigl( z_{\alpha}(s); x(s) \bigr) 
        &= a \angle \bigl( z_{\alpha}(s), y_{0} \bigr) 
          + \ell_{\alpha} \angle \bigl( z_{\alpha}(s), z_{\alpha}(s) \bigr)  
        + \ell_{-\alpha} \angle \bigl( z_{\alpha}(s), -z_{\alpha} \bigr) \notag \\
          & \qquad + \sum_{\gamma \text{ early}} \ell_{\gamma} 
            \angle \bigl(  z_{\gamma}, z_{\alpha}(s) \bigr) 
                + \sum_{\gamma \text{ late}} \ell_{\gamma} 
                  \angle \bigl( z_{\gamma}, z_{\alpha}(s) \bigr) \notag \\
        &= a \angle(z_{\alpha}, y_{0}) + a \Delta_{0}(s) \, s 
          + \ell_{\alpha} \angle( z_{\alpha}, z_{\alpha} )
        + \ell_{-\alpha} \angle(z_{\alpha}, -z_{\alpha}) - \ell_{-\alpha} |s|  \\
          & \qquad + \sum_{\gamma \text{ early}} \ell_{\gamma} 
            \bigl[ \angle( z_{\gamma}, z_{\alpha} ) + s \bigr]
                + \sum_{\gamma \text{ late}} \ell_{\gamma} 
                  \bigl[ \angle( z_{\gamma}, z_{\alpha} ) - s \bigr] \notag \\
        &= G_{a}(z_{\alpha}; x) 
          + \left( a \Delta_{0}(s) - sign(s) \, \ell_{-\alpha} 
            + \sum_{\gamma \text{ early}} \ell_{\gamma} 
              - \sum_{\gamma \text{ late}} \ell_{\gamma} \right) \notag s \\
        &= g + \left( a \Delta_{0}(s) - sign(s) \, \ell_{-\alpha} 
            + \sum_{\gamma \text{ early}} \ell_{\gamma} 
              - \sum_{\gamma \text{ late}} \ell_{\gamma} \right) s \notag .
    \end{align}
since $z_{\alpha} \in M_{a}(x)$, by assumption.

Suppose
     \begin{equation*}
       z_{\alpha} = - y_{0} \in M_{a}(x) .
    \end{equation*}
Choose $w \in S^{1}$ s.t.\ $w \perp z_{\alpha}$ and the sign of $z_{\gamma} \cdot w$ is positive for late $\gamma$ and negative for early. Since $z_{\alpha} \in M_{a}(x)$, we must have 
$\tfrac{d_{-}}{dv} G_{a}(v;x) \restriction_{v=z_{\alpha}} < 0$ and 
$\tfrac{d_{+}}{dv} G_{a}(v;x)  \restriction_{v=z_{\alpha}} > 0$. Thus, applying \eqref{E:angle.deriv.sign} and \eqref{E:one-sided.angle.derivs} then moving the one-sided derivative at $z_{\alpha}$ over to the RHS we get:
    \begin{align}  \label{E:LEFT.RIGHT.ineqs.z.alpha.=.-y0}
       \emph{LEFT: } &a + \ell_{-\alpha} 
         + \sum_{\gamma \text{ early}} \ell_{\gamma} 
                  - \sum_{\gamma \text{ late}} \ell_{\gamma} < \ell_{\alpha} , \\
       \emph{RIGHT: }  -&a - \ell_{-\alpha} 
         + \sum_{\gamma \text{ early}} \ell_{\gamma} 
                  - \sum_{\gamma \text{ late}} \ell_{\gamma} > - \ell_{\alpha} \notag .
    \end{align}

By \eqref{E:angle(y0,z.alpha(s))}, $\Delta_{0}(s)= - sign(s)$. Therefore, the LHS of the first of these inequalities is the coefficient 
of $s$ in \eqref{E:Ga(z.alpha(s),x(s)} when $s < 0$. The RHS the first inequality is the coefficient of $s$ in \eqref{E:Ga(z.beta(s),x(s)} corresponding to early $\beta$. Multiplying both sides 
by $s < 0$ we get, by \eqref{E:Ga(z.alpha(s),x(s)} and \eqref{E:Ga(z.beta(s),x(s)},
    \begin{equation*}
      G_{a} \bigl( z_{\alpha}(s); x(s) \bigr) 
        = g + \left( a + \ell_{-\alpha} + \sum_{\gamma \text{ early}} \ell_{\gamma} 
          - \sum_{\gamma \text{ late}} \ell_{\gamma} \right) s
            > g + \ell_{\alpha} s = G_{a} \bigl( z_{\beta}(s); x(s) \bigr),
    \end{equation*}
providing $\beta$ is early and $s < 0$ is close to 0. 
Thus, $z_{\alpha}(s) \notin M_{a} \bigl( x(s) \bigr)$

The LHS of the second inequality in \eqref{E:LEFT.RIGHT.ineqs.z.alpha.=.-y0}, is the coefficient of $s$ in \eqref{E:Ga(z.alpha(s),x(s)} when $s > 0$. The RHS the second inequality is the coefficient of $s$ in \eqref{E:Ga(z.beta(s),x(s)} corresponding to late $\beta$. Proceeding as in the last paragraph we again find $z_{\alpha}(s) \notin M_{a} \bigl( x(s) \bigr)$.

Similarly, suppose 
    \begin{equation*}
       z_{\alpha} = y_{0} \in M_{a}(x) .
    \end{equation*} 
Let $w \in S^{1}$ be as before. Since $z_{\alpha} \in M_{a}(x)$, we must have $\tfrac{d_{-}}{dv} \restriction_{v=z_{\alpha}} < 0$ and 
$\tfrac{d_{+}}{dv} \restriction_{v=z_{\alpha}} > 0$. Thus, applying \eqref{E:angle.deriv.sign} and \eqref{E:one-sided.angle.derivs} we get:
    \begin{align*}
       \emph{LEFT: } -&a + \ell_{-\alpha} 
         + \sum_{\gamma \text{ early}} \ell_{\gamma} 
                  - \sum_{\gamma \text{ late}} \ell_{\gamma} < \ell_{\alpha} , \\
       \emph{RIGHT: }  &a - \ell_{-\alpha} 
         + \sum_{\gamma \text{ early}} \ell_{\gamma} 
                  - \sum_{\gamma \text{ late}} \ell_{\gamma} > - \ell_{\alpha} .
    \end{align*}

By \eqref{E:angle(y0,z.alpha(s))}, $\Delta_{0}(s)= sign(s)$. Therefore, the LHS of the first of these inequalities is the coefficient of $s$ in \eqref{E:Ga(z.alpha(s),x(s)} when $s < 0$. The RHS of the first inequality is the coefficient of $s$ in \eqref{E:Ga(z.beta(s),x(s)} corresponding to early $\beta$. 
The LHS of the second inequality is the coefficient of $s$ in \eqref{E:Ga(z.alpha(s),x(s)} when $s > 0$. The RHS the second inequality is the coefficient of $s$ in \eqref{E:Ga(z.beta(s),x(s)} corresponding to late $\beta$. Arguing as before we again conclude that 
$z_{\alpha}(s) \notin M_{a} \bigl( x(s) \bigr)$.

Finally, consider the case 
    \begin{equation*}
      z_{\alpha} \in M_{a}(x) \setminus \{ \pm y_{0} \} .
    \end{equation*} 
Then, from \eqref{E:angle(y0,z.alpha(s))}, we have 
$\Delta_{0}(s) = \bigl( 2 E(0) - 1 \bigr)$. Hence, the usual argument leads to the inequalities
    \begin{align*}
       \emph{LEFT: } &a \bigl( 2 E(0) - 1 \bigr) + \ell_{-\alpha} 
         + \sum_{\gamma \text{ early}} \ell_{\gamma} 
           - \sum_{\gamma \text{ late}} \ell_{\gamma} < \ell_{\alpha} , \\
       \emph{RIGHT: }  &a \bigl( 2 E(0) - 1 \bigr) - \ell_{-\alpha} 
          + \sum_{\gamma \text{ early}} \ell_{\gamma} 
            - \sum_{\gamma \text{ late}} \ell_{\gamma} > - \ell_{\alpha} .
    \end{align*}
By the usual argument and the results in the previous two cases, we conclude that when $z_{\alpha} \notin \{ \pm y_{0} \}$ and $s \neq 0$ but close to 0 then $z_{\alpha}(s) \notin M_{a} \bigl( x(s) \bigr)$. 
 
We conclude that by perturbing $z_{\alpha}$ by $s$ close to zero we get $z_{\alpha}(s) \notin M_{a} \bigl( x(s) \bigr)$. Hence, \eqref{E:precisely.in.Ma(x(s))} holds in general.

We observed above that $\beta = 1$ is early and $\beta = i$ is late. Hence, after the perturbation, exactly one of $z_{1}$ and $z_{i}$ is in $M_{a} \bigl( x(s) \bigr)$. In particular, $0 < \bigl| M_{a} \bigl( x(s) \bigr) \bigr| < \bigl| M_{a}(x) \bigr|$, 
where $|S|$ is the cardinality of a set $S$. But the important point is that with $|s|$ small, we get data sets $x_{-}(s)$ and $x_{+}(s)$ corresponding to $s<0$ and $s>0$, resp., s.t.\ $M_{a} \bigl( x_{+}(s) \bigr)$ and $M_{a} \bigl( x_{-}(s) \bigr)$ 
are nonempty disjoint subsets of, resp., the late and early sets of points in $x$. 
Importantly, $x_{\pm}(s) \to x$ as $s \to 0$. 

Now, apply this procedure recursively to each of $x_{-}(s)$ and $x_{+}(s)$. We eventually end up with at least two data sets $x_{minus}(\mbf{s})$ and 
$x_{plus}(\mbf{s}')$ depending on vectors $\mbf{s}$ and $\mbf{s}'$ close to 0 and having the following properties. First, $M_{a}\bigl( x_{minus}(\mbf{s}) \bigr)$ contains exactly one point in $S^{1}$, one of the early points in $M_{a}(x)$ and $M_{a}\bigl( x_{plus}(\mbf{s}) \bigr)$ contains exactly one point, one of the late points in $M_{a}(x)$. This means 
$x_{minus}(\mbf{s}), x_{plus}(\mbf{s}') \in \D'$. Second, 
as $\mbf{s}, \mbf{s}', s \to 0$ we have 
$x_{minus}(\mbf{s}), x_{plus}(\mbf{s}') \to x$, our original data set 
in $\D \setminus \D'$. This means $x$ is in the closure of $\D'$. 
Since $x \in \D \setminus \D'$ is arbitrary this shows $\D'$ is dense in $\D$. 

The singleton $M_{a}\bigl( x_{minus}(\mbf{s}) \bigr)$ is among the early points of $M_{a}(x)$ and $M_{a}\bigl( x_{plus}(\mbf{s}) \bigr)$ is among the late. The early and late sets are a positive distance from each other. It follows that $x$ is a singularity (w.r.t.\ $\D'$). 
  \end{proof}

\end{document}
